\documentclass{../yalumba}

\title{From Non-Determinism to Canonical Extraction:\\[0.3em]
Diagnosing and Fixing the Representation Bottleneck\\[0.2em]
in Structural Genomics}

\author{%
  Ajax Davis\\[0.3em]
  \normalsize\textit{Yalumba Project}\\[0.2em]
  \normalsize\texttt{github.com/ajaxdavis/yalumba}%
}

\date{April 2026}

\begin{document}
\maketitle
\thispagestyle{empty}

% ============================================================================
% ABSTRACT
% ============================================================================
\begin{abstract}
\noindent
The yalumba project attempts to infer inherited genomic organization---multi-scale
``ecosystem'' signatures---directly from raw sequencing reads, without alignment
or variant calling. Over ten iterative versions of the Spectral Ecology algorithm,
we deployed increasingly sophisticated mathematical machinery: Laplacian
eigendecomposition, heat kernel signatures, Ollivier--Ricci curvature, diffusion
distances, optimal transport, persistence fields, and Sinkhorn-symmetrized kernels.
Despite this, the best gap between related and unrelated pairs reached only
$-0.61\%$, and zero of fifteen extracted features achieved a signal-to-nuisance
ratio (SNR) above the threshold of~3.

A systematic representation audit---comprising self-perturbation testing under
read shuffling and a coverage ladder analysis---revealed the root cause: the
module extraction layer was \textbf{non-deterministic}. The same sample processed
with the same reads in a different order produced wildly different module graphs.
Module counts changed by 77--105\% under shuffling. Triangle counts changed by
up to 607\%. Every downstream feature, no matter how theoretically invariant,
was built on sand.

The non-determinism traced to three interacting bugs in the module builder:
(1)~a fixed-size hash table with early termination at 70\% load factor, causing
different read orderings to fill different co-occurrence pairs before the table
saturated; (2)~an integer pair hash with collision pathology; and (3)~unsorted
BFS traversal producing order-dependent component labeling. The fix---a canonical
builder using complete pair enumeration, frequency-threshold cutoff, collision-free
string keys, and sorted BFS---achieves perfect determinism: identical module counts,
identical pair counts (6,016,446), and identical graph structure across all tested
shuffles.

This paper documents the full journey: the ten-version spectral ecology arc that
motivated the audit, the audit protocol and its devastating results, the line-by-line
root cause analysis, the canonical builder design, and the determinism verification.
We argue that this is the project's most important single finding. The spectral
ecology program was not wrong in theory---it was measuring noise masquerading as
structure. The canonical builder provides, for the first time, a deterministic
substrate on which structural genomics methods can be honestly evaluated.

The result is a cautionary tale for computational biology: sophisticated
downstream analysis cannot rescue a non-deterministic upstream representation.
No amount of spectral geometry, optimal transport, or topological persistence
can extract signal from a feature space dominated by hash table insertion
artifacts. The path forward requires measurement science discipline---acceptance
gates on representation stability before any kinship model is attempted.
\end{abstract}

\newpage
\tableofcontents
\newpage

% ============================================================================
% 1. INTRODUCTION
% ============================================================================
\section{Introduction}
\label{sec:introduction}

The yalumba project's central hypothesis is that human inheritance can be modeled
as the transmission of resilient, interacting genomic ecosystems whose structure
can be inferred from raw sequencing reads without reference coordinates or variant
calling~\cite{thompson2013identity}. Relatedness between two individuals is not
primarily a matter of shared genomic tokens---shared k-mers, shared haplotype
segments---but of shared \textit{organizational principles}: do these two genomes
instantiate the same \textit{kind} of structure?

This hypothesis led to the Spectral Ecology program: a sequence of algorithms
that extract module co-occurrence graphs from k-mer data, compute spectral
decompositions of the resulting graph Laplacians, and compare the spectral
signatures across samples to infer kinship. The mathematical framework draws on
spectral graph theory~\cite{chung1997spectral}, heat kernel
signatures~\cite{sun2009concise}, diffusion maps~\cite{coifman2006diffusion},
discrete Ricci curvature~\cite{jost2011riemannian}, and optimal
transport~\cite{villani2009optimal}.

\subsection{The Ten-Version Paradox}

Over ten iterative versions (v1--v10), the Spectral Ecology algorithm accumulated
layer upon layer of mathematical sophistication:

\begin{itemize}[leftmargin=2em]
  \item \textbf{v1--v2:} Dense co-occurrence graphs, Laplacian eigendecomposition,
    heat kernel signatures, basic ecological role assignment.
  \item \textbf{v3--v5:} Optimal transport (Wasserstein distance) on spectral
    distributions, coherence fields, multi-scale analysis.
  \item \textbf{v6:} Persistence fields via topological filtration of the module graph.
  \item \textbf{v7:} Sinkhorn-symmetrized transport kernels, achieving the best gap
    of $-0.61\%$ with 55\% cooperative classification.
  \item \textbf{v8--v10:} Curvature-weighted features, coverage deconfounding layers,
    feature selection by stability ranking.
\end{itemize}

The gap trajectory told a story of gradual improvement: $-26.87\% \to -3.21\% \to
-0.61\%$. But it never crossed zero. The two algorithms that \textit{did} achieve
positive gaps---Module Persistence v3 (+3.36\% separation, +0.13\% gap) and
Coalition Transfer v3 (+1.97\% separation, +0.06\% gap)---used comparatively
simple features. The most mathematically ambitious algorithm, consuming the
richest feature space, performed worst.

This was the paradox: \textbf{increasing downstream sophistication produced
decreasing returns}. The natural interpretation was that the invariants were
wrong, or the theory was wrong, or genomic ecosystems simply didn't carry enough
kinship signal. An external review suggested a different hypothesis: perhaps the
\textit{representation itself} was the problem.

\subsection{The External Feedback}

The critical insight came from outside the project: before adding more mathematical
machinery on top, we should verify that the upstream representation---the module
graph---was stable enough to support \textit{any} downstream analysis. This meant
asking two questions that had never been systematically tested:

\begin{enumerate}[leftmargin=2em]
  \item \textbf{Self-perturbation:} If we process the same sample's reads in a
    different order, do we get the same module graph?
  \item \textbf{Coverage sensitivity:} If we subsample reads from the same sample,
    how much do the features change?
\end{enumerate}

The answers were catastrophic.

\subsection{Why Negative Results Are the Most Important Results}

This paper documents a negative result: the spectral ecology program, in its
ten-version incarnation, was built on a non-deterministic foundation. None of
its results---positive or negative---can be trusted as measurements of the
underlying biology.

We argue this is the project's most important finding to date. It is more
valuable than any positive kinship gap because it identifies a \textit{systematic}
error that would have corrupted every future algorithm built on the same
representation. Finding and fixing it converts the entire module extraction
layer from a source of uncontrolled noise into a calibrated instrument.

The structure of the paper follows the diagnostic arc: the audit that exposed
the problem (Section~\ref{sec:audit}), the root cause analysis that explained it
(Section~\ref{sec:rootcause}), the canonical builder that fixes it
(Section~\ref{sec:canonical}), the verification that the fix works
(Section~\ref{sec:verification}), the reinterpretation of all prior results
(Section~\ref{sec:reinterpretation}), and the path forward
(Section~\ref{sec:forward}).


% ============================================================================
% 2. BACKGROUND
% ============================================================================
\section{Background: The Module Extraction Pipeline}
\label{sec:background}

Before presenting the audit results, we briefly describe the module extraction
pipeline whose stability is under investigation.

\subsection{From Reads to K-mer Co-occurrence}

Given a set of sequencing reads for a sample, the pipeline:
\begin{enumerate}[leftmargin=2em]
  \item Extracts all canonical k-mers from each read using a rolling hash
    (Rabin--Karp polynomial, Mersenne prime modulus $2^{61}-1$).
  \item For each read, enumerates all pairs of k-mers that co-occur within
    the same read.
  \item Accumulates co-occurrence counts into a pair frequency table.
  \item Filters pairs by frequency to retain only statistically significant
    co-occurrences.
  \item Constructs a co-occurrence graph where nodes are k-mers and edges
    represent retained co-occurrences.
\end{enumerate}

\subsection{From Co-occurrence Graph to Modules}

The co-occurrence graph is decomposed into connected components, each of which
is called a \textit{module}. A module represents a set of k-mers that
co-occur frequently enough to form a coherent unit---analogous to a ``motif
society'' in the project's ecological framework.

\subsection{From Modules to Features}

The module graph (nodes = modules, edges = inter-module co-occurrences) is then
characterized by a battery of features:
\begin{itemize}[leftmargin=2em]
  \item \textbf{Structural:} module count, edge count, edge density, triangle count.
  \item \textbf{Degree:} mean degree, max degree.
  \item \textbf{Curvature:} Ollivier--Ricci curvature statistics (mean, std, P50, P90).
  \item \textbf{Ecological roles:} fraction of nodes classified as core, satellite,
    bridge, scaffold, or boundary based on degree and curvature profiles.
\end{itemize}

It is this feature vector that the Spectral Ecology algorithm uses for kinship
comparison. The stability of this vector is the subject of the audit.


% ============================================================================
% 3. THE REPRESENTATION AUDIT
% ============================================================================
\section{The Representation Audit}
\label{sec:audit}

\subsection{Experimental Protocol}

The audit was designed to answer one question: \textit{is the module extraction
pipeline a calibrated measurement instrument?} A calibrated instrument produces
the same measurement when applied to the same sample under trivial perturbations.
Two perturbation classes were tested:

\begin{definition}[Self-Perturbation (Shuffle Test)]
Given a sample's reads $R = \{r_1, r_2, \ldots, r_n\}$, construct shuffled
versions $R_\sigma = \{r_{\sigma(1)}, r_{\sigma(2)}, \ldots, r_{\sigma(n)}\}$
for random permutations $\sigma$. Run the full extraction pipeline on both $R$
and $R_\sigma$. Compare the resulting feature vectors.

\textit{A deterministic pipeline must produce identical outputs.}
\end{definition}

\begin{definition}[Coverage Ladder]
Given a sample's reads $R$, construct subsampled versions at fractions
$f \in \{0.25, 0.50, 0.75, 1.0\}$. Run the pipeline on each. Measure how
features scale with coverage.

\textit{A coverage-robust pipeline should show smooth, predictable scaling.}
\end{definition}

Three random seeds were used for shuffling: 42, 12345, and 99999. Two samples
from the CEPH~1463 pedigree~\cite{eberle2017platinum} were tested: NA12889
(grandfather) and a second family member. Each configuration was run through
the full pipeline with 50,000 reads per sample.

\subsection{Feature Extraction}
\label{sec:features}

Fifteen features were extracted from each module graph, spanning four categories:

\begin{table}[H]
\centering
\caption{Feature catalogue: 15 features across 4 categories.}
\label{tab:features}
\begin{tabular}{lll}
\toprule
\rowcolor{table-header}
\textcolor{white}{\textbf{Category}} & \textcolor{white}{\textbf{Feature}} & \textcolor{white}{\textbf{Description}} \\
\midrule
Structural & \texttt{moduleCount}    & Number of connected components \\
           & \texttt{triangleCount}  & Number of triangles in module graph \\
           & \texttt{edgeCount}      & Number of edges in module graph \\
           & \texttt{edgeDensity}    & Edge count / max possible edges \\
\midrule
Degree     & \texttt{meanDegree}     & Average node degree \\
           & \texttt{maxDegree}      & Maximum node degree \\
\midrule
Curvature  & \texttt{curvMean}       & Mean Ollivier--Ricci curvature \\
           & \texttt{curvStd}        & Std.\ dev.\ of curvature \\
           & \texttt{curvP50}        & Median curvature \\
           & \texttt{curvP90}        & 90th percentile curvature \\
\midrule
Ecological & \texttt{role\_core}      & Fraction of core nodes \\
           & \texttt{role\_satellite} & Fraction of satellite nodes \\
           & \texttt{role\_bridge}    & Fraction of bridge nodes \\
           & \texttt{role\_scaffold}  & Fraction of scaffold nodes \\
           & \texttt{role\_boundary}  & Fraction of boundary nodes \\
\bottomrule
\end{tabular}
\end{table}

\subsection{Self-Perturbation Analysis}

The self-perturbation test revealed the core problem. Under the original builder,
shuffling the reads of sample NA12889 produced dramatically different module
graphs:

\begin{table}[H]
\centering
\caption{Module count under read shuffling (original builder, NA12889).}
\label{tab:shuffle-modules}
\begin{tabular}{lrrr}
\toprule
\rowcolor{table-header}
\textcolor{white}{\textbf{Ordering}} & \textcolor{white}{\textbf{Modules}} & \textcolor{white}{$\Delta$} & \textcolor{white}{\textbf{Change}} \\
\midrule
Original          & 457 & ---   & --- \\
Shuffle (seed 42)    & 938 & +481 & +105\% \\
Shuffle (seed 12345) & 809 & +352 & +77\% \\
Shuffle (seed 99999) & 818 & +361 & +79\% \\
\bottomrule
\end{tabular}
\end{table}

\begin{keyfinding}[Shuffle Catastrophe]
The same sample, with the same reads, processed in a different order, produces
anywhere from 457 to 938 modules---a variation of 77--105\%. \textbf{The module
extraction pipeline is not a measurement instrument; it is a random number
generator seeded by read ordering.}
\end{keyfinding}

The instability was not confined to module count. The broader shuffle test
across all samples showed:

\begin{itemize}[leftmargin=2em]
  \item Module count changes of 45--184\%.
  \item Triangle count changes of up to 607\%.
  \item Edge count and density variations proportional to module count changes.
  \item Curvature statistics fluctuating wildly, with self-variation exceeding
    cross-sample variation.
\end{itemize}

\subsection{Cross-Sample Variance}

For a feature to be useful for kinship discrimination, it must vary \textit{more}
between unrelated samples than within repeated measurements of the same sample.
We quantified this via the coefficient of variation (CV):

\begin{equation}
  \text{CV}_{\text{self}} = \frac{\text{std}(\text{feature across shuffles of same sample})}
    {\text{mean}(\text{feature across shuffles of same sample})}
\end{equation}

\begin{equation}
  \text{CV}_{\text{cross}} = \frac{\text{std}(\text{feature across different samples})}
    {\text{mean}(\text{feature across different samples})}
\end{equation}

For a feature to carry kinship signal, we require $\text{CV}_{\text{cross}} \gg
\text{CV}_{\text{self}}$.

\subsection{Signal-to-Nuisance Ratio}
\label{sec:snr}

We define the signal-to-nuisance ratio (SNR) as:

\begin{equation}
  \text{SNR} = \frac{\text{CV}_{\text{cross}}}{\text{CV}_{\text{self}}}
  \label{eq:snr}
\end{equation}

An SNR of 1 means the feature varies as much within a sample (due to processing
artifacts) as between samples (due to biology). An SNR above 3 is the minimum
threshold for a feature to be considered a candidate for kinship discrimination.
Table~\ref{tab:snr} presents the full results.

\begin{table}[H]
\centering
\caption{Signal-to-nuisance ratio for all 15 features. Threshold: SNR $> 3$.
\textbf{Zero features pass.}}
\label{tab:snr}
\begin{tabular}{lrrrp{5.5cm}}
\toprule
\rowcolor{table-header}
\textcolor{white}{\textbf{Feature}} & \textcolor{white}{\textbf{Self CV}} & \textcolor{white}{\textbf{Cross CV}} & \textcolor{white}{\textbf{SNR}} & \textcolor{white}{\textbf{Verdict}} \\
\midrule
\texttt{moduleCount}      & 0.1059 & 0.1536 & 1.45 & \textcolor{yalumba-gold}{MARGINAL} \\
\texttt{triangleCount}    & 0.4386 & 0.4125 & 0.94 & \textcolor{red}{NOISE $>$ SIGNAL} \\
\texttt{edgeCount}        & 0.2303 & 0.2353 & 1.02 & \textcolor{yalumba-gold}{MARGINAL} \\
\texttt{edgeDensity}      & 0.0725 & 0.2124 & 2.93 & \textcolor{yalumba-gold}{MARGINAL (best)} \\
\texttt{meanDegree}       & 0.1249 & 0.1465 & 1.17 & \textcolor{yalumba-gold}{MARGINAL} \\
\texttt{maxDegree}        & 0.2947 & 0.4739 & 1.61 & \textcolor{yalumba-gold}{MARGINAL} \\
\texttt{curvMean}         & 0.1373 & 0.1244 & 0.91 & \textcolor{red}{NOISE $>$ SIGNAL} \\
\texttt{curvStd}          & 0.2452 & 0.2275 & 0.93 & \textcolor{red}{NOISE $>$ SIGNAL} \\
\texttt{curvP50}          & 0.0660 & 0.1633 & 2.47 & \textcolor{yalumba-gold}{MARGINAL} \\
\texttt{curvP90}          & 0.1949 & 0.2717 & 1.39 & \textcolor{yalumba-gold}{MARGINAL} \\
\texttt{role\_core}       & 0.5458 & 0.5315 & 0.97 & \textcolor{red}{NOISE $>$ SIGNAL} \\
\texttt{role\_satellite}  & 0.4689 & 0.7827 & 1.67 & \textcolor{yalumba-gold}{MARGINAL} \\
\texttt{role\_bridge}     & 1.0875 & 2.2302 & 2.05 & \textcolor{yalumba-gold}{MARGINAL} \\
\texttt{role\_scaffold}   & 0.8034 & 0.4767 & 0.59 & \textcolor{red}{NOISE $>$ SIGNAL} \\
\texttt{role\_boundary}   & 0.7251 & 0.8576 & 1.18 & \textcolor{yalumba-gold}{MARGINAL} \\
\bottomrule
\end{tabular}
\end{table}

\begin{keyfinding}[Zero Features Pass]
Of 15 extracted features, \textbf{zero} achieve SNR $> 3$. The best
feature---\texttt{edgeDensity} at 2.93---is still below threshold. Five features
have SNR $< 1$, meaning their self-perturbation variance \textit{exceeds} their
cross-sample variance. The instrument is noisier than the signal it is trying
to measure.
\end{keyfinding}

\subsection{The Shuffle Catastrophe in Detail}

The magnitude of the shuffle effect deserves emphasis. Consider the curvature
features, which were the theoretical centerpiece of the spectral ecology program.
Ollivier--Ricci curvature was chosen precisely because it is a \textit{geometric
invariant}---it should capture intrinsic structure independent of labeling or
ordering. Yet:

\begin{itemize}[leftmargin=2em]
  \item \texttt{curvMean}: SNR = 0.91. The curvature mean varies \textit{more}
    across shuffles of the same sample than across genuinely different samples.
  \item \texttt{curvStd}: SNR = 0.93. Same pathology.
  \item The curvature distribution is not reflecting graph geometry; it is
    reflecting hash table insertion order.
\end{itemize}

The triangle count showed the most extreme instability: changes of up to 607\%
under shuffling. Since the number of triangles in a graph is a cubic function of
the adjacency structure, even small changes in which edges are present cause
large changes in triangle count. The triangle count was amplifying the
non-determinism rather than averaging it out.

\subsection{Interpretation: Downstream Sophistication Outrunning Upstream
Identifiability}

The audit results admit a clean interpretation. The module extraction pipeline
has two layers:

\begin{enumerate}[leftmargin=2em]
  \item \textbf{Upstream:} Read processing $\to$ pair counting $\to$ graph
    construction $\to$ module extraction. This layer defines the representation.
  \item \textbf{Downstream:} Feature extraction $\to$ spectral analysis $\to$
    curvature computation $\to$ kinship scoring. This layer analyzes the
    representation.
\end{enumerate}

The ten versions of Spectral Ecology invested exclusively in the downstream
layer. Each version added more sophisticated invariants, more careful normalization,
more deconfounding. But the upstream layer---the representation itself---was never
validated.

\begin{limitation}[The Fundamental Error]
We spent ten versions optimizing the analysis of a measurement that was not
reproducible. This is the definition of a calibration failure. No amount of
downstream sophistication can rescue a non-deterministic upstream representation.
\end{limitation}

The analogy to experimental science is precise: it is as if we built increasingly
sensitive detectors to measure a quantity, while the measurement apparatus itself
had an uncontrolled source of noise larger than the signal. Better detectors
would detect the noise more precisely, not the signal.


% ============================================================================
% 4. ROOT CAUSE ANALYSIS
% ============================================================================
\section{Root Cause Analysis}
\label{sec:rootcause}

With the non-determinism established, we traced the failure to its source in the
codebase. The root cause was localized to a single file---\texttt{module-builder.ts}---and
three interacting implementation decisions.

\subsection{The Early Termination Bug}
\label{sec:earlytermination}

Line 74 of \texttt{module-builder.ts} contained:

\begin{lstlisting}[language=Java,caption={The early termination bug.}]
if (insertions > MAX_LOAD) break;
\end{lstlisting}

The module builder used a fixed-size hash table with $2^{24} = 16{,}777{,}216$
entries to store k-mer co-occurrence pairs. The \texttt{MAX\_LOAD} constant was
set to 70\% of this capacity---approximately 11.7 million entries. When the
number of unique pair insertions exceeded this threshold, the builder stopped
processing reads entirely.

This is an \textit{early termination} strategy: rather than resizing the hash
table (as a standard hash map would), the builder silently discards all remaining
reads. The rationale was presumably performance---avoid the cost of table
resizing for a streaming pipeline. The consequence was catastrophic.

\subsection{Why a Fixed-Size Hash Table Causes Order Dependence}

The mechanism is straightforward:

\begin{enumerate}[leftmargin=2em]
  \item The hash table has $N = 16\text{M}$ slots.
  \item Reads are processed sequentially, extracting k-mer pairs.
  \item Each new pair is inserted into the table.
  \item When the number of unique pairs exceeds $0.7N \approx 11.7\text{M}$,
    the builder terminates.
  \item The set of pairs in the table depends on which reads were processed
    \textit{before} the capacity limit was hit.
\end{enumerate}

Different read orderings cause different pairs to be inserted before saturation.
The saturated table represents a \textit{random sample} of the true co-occurrence
structure, where the sampling distribution is determined by read ordering.

\begin{definition}[Order-Dependent Saturation]
Let $P$ be the set of all k-mer co-occurrence pairs in a sample, with $|P| > 0.7N$.
A fixed-size hash table of capacity $N$ can store at most $0.7N$ unique pairs
before early termination. The stored subset $P_\sigma \subset P$ depends on the
read permutation $\sigma$:
\begin{equation}
  P_\sigma \neq P_{\sigma'} \quad \text{for } \sigma \neq \sigma' \text{ (in general)}
\end{equation}
The module graph $G_\sigma$ constructed from $P_\sigma$ is therefore a function
of $\sigma$, not just of the sample's biology.
\end{definition}

For sample NA12889 with 50,000 reads, the canonical builder (which processes all
reads) found 6,016,446 unique co-occurrence pairs. The original builder's hash
table could store at most $\sim$11.7M pairs, which would seem sufficient---but the
actual behavior depends on the hash function's collision rate. With collisions,
the effective capacity was lower, and many samples exceeded it.

\subsection{Hash Collisions in Pair Counting}

The pair hash function used integer arithmetic:

\begin{lstlisting}[language=Java,caption={The collision-prone pair hash.}]
const hash = ((a * 0x9e3779b9) ^ b) >>> 0;
\end{lstlisting}

Here \texttt{a} and \texttt{b} are k-mer hash values, and \texttt{0x9e3779b9}
is the integer part of $\phi^{-1} \cdot 2^{32}$ (the golden ratio constant
commonly used in hash functions). This hash has two problems:

\begin{enumerate}[leftmargin=2em]
  \item \textbf{Non-commutativity:} $h(a, b) \neq h(b, a)$ in general, but
    k-mer co-occurrence is symmetric. If the code does not consistently canonicalize
    pair ordering, the same biological co-occurrence may hash to different slots
    depending on which k-mer is encountered first.
  \item \textbf{Collision rate:} The 32-bit hash space ($2^{32} \approx 4.3\text{B}$
    values) mapped into a $2^{24}$ table via modular reduction. With millions of
    pairs, the birthday paradox guarantees significant collisions. Each collision
    means a pair that \textit{should} be counted is instead lost or conflated
    with another pair.
\end{enumerate}

Collisions interact with early termination: colliding pairs reduce the effective
capacity, causing the table to saturate earlier, causing more reads to be
discarded, amplifying the order dependence.

\subsection{Non-Deterministic BFS Traversal}

After pair counting, the module builder constructs an adjacency list and performs
breadth-first search (BFS) to find connected components. The BFS used JavaScript's
\texttt{Map} iteration order, which is insertion order. Since the adjacency list
was built from the hash table (whose contents depend on read ordering), the BFS
visited nodes in a data-dependent order, producing different component labelings
for different runs.

This is a subtler source of non-determinism than the pair counting issue. Even
if two runs produced identical adjacency structures, different BFS traversal
orders could produce different component boundaries if the graph contained
near-disconnected regions (components connected by a single edge that may or
may not be present depending on hash collisions).

\subsection{The Cascade}
\label{sec:cascade}

The three sources of non-determinism---early termination, hash collisions, and
BFS ordering---compound through the pipeline:

\begin{equation}
\boxed{
\begin{aligned}
  &\text{Read ordering} \\
  &\quad \xrightarrow{\text{early termination}} \text{Different pairs counted} \\
  &\quad \xrightarrow{\text{hash collisions}} \text{Different adjacency structure} \\
  &\quad \xrightarrow{\text{BFS ordering}} \text{Different components} \\
  &\quad \xrightarrow{\text{module labeling}} \text{Different modules} \\
  &\quad \xrightarrow{\text{feature extraction}} \text{Different feature vectors} \\
  &\quad \xrightarrow{\text{comparison}} \text{Different kinship scores}
\end{aligned}
}
\label{eq:cascade}
\end{equation}

Each stage amplifies the non-determinism of the previous stage. The triangle
count's 607\% variation is a consequence of this amplification: triangles require
three mutually adjacent nodes, so even a small change in edge set produces a
large change in triangle count (the change is cubic in the edge perturbation
for dense subgraphs).

\begin{keyfinding}[Root Cause Identified]
The non-determinism originates at line 74 of \texttt{module-builder.ts}:
\texttt{if (insertions > MAX\_LOAD) break}. A fixed-size hash table with early
termination causes the set of counted co-occurrence pairs to depend on read
ordering. Hash collisions and unsorted BFS amplify this into order-dependent
module extraction. The downstream spectral ecology features are functions of
this random process, not of the underlying biology.
\end{keyfinding}


% ============================================================================
% 5. THE CANONICAL BUILDER
% ============================================================================
\section{The Canonical Builder}
\label{sec:canonical}

The canonical module builder was designed from first principles to eliminate
every source of non-determinism identified in Section~\ref{sec:rootcause}.

\subsection{Design Principles}

Three requirements governed the design:

\begin{enumerate}[leftmargin=2em]
  \item \textbf{Determinism:} The output must be a pure function of the input
    \textit{set} of reads, invariant to ordering. Same reads $\Rightarrow$ same
    modules, regardless of permutation.
  \item \textbf{Completeness:} All reads must be processed. No early termination,
    no capacity limits, no silent data loss.
  \item \textbf{Collision freedom:} The pair counting data structure must not
    conflate distinct pairs. Every unique co-occurrence must be counted exactly
    once.
\end{enumerate}

These requirements trade performance for correctness. The canonical builder is
slower than the original (177--492 seconds per extraction versus the original's
faster but incorrect execution). This tradeoff is acceptable: correctness is a
prerequisite, and performance can be recovered later via native code
acceleration (Section~\ref{sec:spirv}).

\subsection{Fix 1: No Early Termination}

The canonical builder uses a dynamically-sized data structure (JavaScript
\texttt{Map}) instead of a fixed-size hash table. There is no capacity limit
and no early termination. Every read is processed, and every k-mer pair is
counted.

\begin{lstlisting}[language=Java,caption={Canonical builder: no capacity limit.}]
// OLD: fixed-size table with early termination
// if (insertions > MAX_LOAD) break;

// NEW: dynamic map, all reads processed
for (const read of reads) {
  const kmers = extractKmers(read);
  for (const [a, b] of allPairs(kmers)) {
    const key = canonicalKey(a, b);
    pairCounts.set(key, (pairCounts.get(key) ?? 0) + 1);
  }
}
// No break. Every read processed.
\end{lstlisting}

\subsection{Fix 2: Frequency-Threshold Cutoff}
\label{sec:threshold}

The original builder used a ``top-$N$'' strategy: retain the $N$ most frequent
pairs. This introduces non-determinism when multiple pairs share the same
frequency at the threshold boundary---the set of retained pairs depends on
which pairs happen to be in which positions in the hash table.

The canonical builder uses a \textit{frequency threshold} instead: retain all
pairs whose count exceeds a minimum frequency $f_{\min}$. This is deterministic
because the frequency of each pair is a deterministic function of the read
\textit{set} (not ordering), and the threshold comparison $\text{count}(p) \geq
f_{\min}$ has no tie-breaking ambiguity.

\begin{equation}
  E_{\text{canonical}} = \{ (a, b) : \text{count}(a, b) \geq f_{\min} \}
  \label{eq:threshold}
\end{equation}

The threshold $f_{\min}$ can be set based on statistical significance (e.g.,
pairs occurring more than expected under a null model of independent k-mer
occurrence) without introducing order dependence.

\subsection{Fix 3: Collision-Free String Keys}

The integer pair hash \texttt{((a * 0x9e3779b9) \^{} b) >>> 0} was replaced
with string keys of the form \texttt{"a:b"} where \texttt{a} and \texttt{b}
are the k-mer hash values in canonical (sorted) order:

\begin{lstlisting}[language=Java,caption={Collision-free string keys.}]
function canonicalKey(a: bigint, b: bigint): string {
  return a < b ? `${a}:${b}` : `${b}:${a}`;
}
\end{lstlisting}

String keys in a JavaScript \texttt{Map} are compared by value, not by hash.
Two distinct pairs can never collide. The cost is higher memory usage and slower
comparison (string comparison vs.\ integer comparison), but the correctness
guarantee is absolute.

The canonical ordering (\texttt{a < b}) ensures that the pair $(a, b)$ and
$(b, a)$ are always stored under the same key, eliminating the non-commutativity
issue of the original hash.

\subsection{Fix 4: Sorted BFS for Deterministic Component Ordering}

The BFS traversal was modified to process neighbors in sorted order:

\begin{lstlisting}[language=Java,caption={Sorted BFS for deterministic components.}]
function bfs(adjacency: Map<number, number[]>, start: number) {
  const queue = [start];
  const visited = new Set([start]);
  const component: number[] = [];

  while (queue.length > 0) {
    const node = queue.shift()!;
    component.push(node);
    // Sort neighbors for deterministic traversal
    const neighbors = adjacency.get(node) ?? [];
    neighbors.sort((a, b) => a - b);
    for (const neighbor of neighbors) {
      if (!visited.has(neighbor)) {
        visited.add(neighbor);
        queue.push(neighbor);
      }
    }
  }
  return component;
}
\end{lstlisting}

With sorted BFS, the component membership and labeling are deterministic
functions of the adjacency structure alone, independent of the order in which
nodes were inserted into the adjacency map.

\subsection{Memory and Performance Tradeoffs}

The canonical builder's design choices have measurable performance costs:

\begin{table}[H]
\centering
\caption{Performance comparison: original vs.\ canonical builder.}
\label{tab:performance}
\begin{tabular}{lll}
\toprule
\rowcolor{table-header}
\textcolor{white}{\textbf{Property}} & \textcolor{white}{\textbf{Original}} & \textcolor{white}{\textbf{Canonical}} \\
\midrule
Pair storage    & Fixed 16M integer slots & Dynamic Map (string keys) \\
Memory          & $\sim$128 MB fixed       & Proportional to unique pairs \\
Early termination & At 70\% capacity       & Never \\
Collision rate  & Non-zero (integer hash)  & Zero (string equality) \\
BFS ordering    & Insertion order          & Sorted (deterministic) \\
Time per sample & Fast but wrong           & 177--492s (correct) \\
\bottomrule
\end{tabular}
\end{table}

\begin{limitation}[Performance Cost]
The canonical builder is 2--5$\times$ slower than the original, primarily due
to string key construction and comparison. This is acceptable for research
correctness but will need to be addressed for production use. The planned
SPIR-V/C acceleration (Section~\ref{sec:spirv}) will recover this performance
using native data structures with deterministic hash functions.
\end{limitation}

\subsection{Future: SPIR-V/C Acceleration}
\label{sec:spirv}

The canonical builder's correctness guarantees can be preserved while recovering
performance via native code:

\begin{enumerate}[leftmargin=2em]
  \item \textbf{C implementation:} Replace JavaScript \texttt{Map} with a
    dynamically-sized open-addressing hash table using a deterministic,
    collision-resistant hash (e.g., xxHash128 on canonically-ordered 128-bit
    pair keys). This eliminates the string overhead while maintaining
    collision freedom for practical purposes.
  \item \textbf{SPIR-V compute shaders:} The pair enumeration step (all pairs
    of k-mers within a read) is embarrassingly parallel. Each read's pairs can
    be computed independently on a GPU workgroup, with atomic accumulation into
    a shared hash table.
  \item \textbf{FFI bridge:} The TypeScript layer calls the native builder via
    Bun's FFI, receiving the completed pair count table and running BFS in
    TypeScript (or also natively).
\end{enumerate}

The key constraint is that the native implementation must produce
\textit{bit-identical} output to the canonical TypeScript builder. This is
verified by the determinism test suite (Section~\ref{sec:verification}).


% ============================================================================
% 6. DETERMINISM VERIFICATION
% ============================================================================
\section{Determinism Verification}
\label{sec:verification}

\subsection{Test Protocol}

The determinism test applies both the original and canonical builders to the
same sample under multiple read orderings:

\begin{enumerate}[leftmargin=2em]
  \item \textbf{Original ordering:} Reads in file order.
  \item \textbf{Shuffle (seed 42):} Reads permuted with PRNG seed 42.
  \item \textbf{Shuffle (seed 12345):} Reads permuted with PRNG seed 12345.
  \item \textbf{Shuffle (seed 99999):} Reads permuted with PRNG seed 99999.
\end{enumerate}

For each configuration, we record: module count, unique pair count, graph node
count, and the full adjacency structure. Two runs are considered a MATCH if and
only if all four quantities are identical.

The test was run on sample NA12889 from the CEPH~1463 pedigree with 50,000
reads. A second sample was being tested at the time of writing.

\subsection{Old Builder Results}

The original builder failed every shuffle test:

\begin{table}[H]
\centering
\caption{Old builder determinism test (NA12889). All shuffles fail.}
\label{tab:old-builder}
\begin{tabular}{lrrrl}
\toprule
\rowcolor{table-header}
\textcolor{white}{\textbf{Ordering}} & \textcolor{white}{\textbf{Modules}} & \textcolor{white}{$\Delta$} & \textcolor{white}{\textbf{Change}} & \textcolor{white}{\textbf{Result}} \\
\midrule
Original           & 457  & ---   & ---      & (baseline) \\
Shuffle (42)       & 938  & +481  & +105\%   & \textcolor{red}{\textbf{FAIL}} \\
Shuffle (12345)    & 809  & +352  & +77\%    & \textcolor{red}{\textbf{FAIL}} \\
Shuffle (99999)    & 818  & +361  & +79\%    & \textcolor{red}{\textbf{FAIL}} \\
\bottomrule
\end{tabular}
\end{table}

Three out of three shuffles produced different module counts. The variation
ranges from 77\% to 105\%. The original builder is \textbf{catastrophically
non-deterministic}.

Note the asymmetry: the original ordering produces \textit{fewer} modules (457)
than any shuffled ordering (809--938). This suggests that the file ordering---which
corresponds to the sequencer's output order---happens to produce a particularly
compact co-occurrence table. Shuffled orderings spread the pairs more evenly
across the hash table, hitting different collision patterns and different
early-termination boundaries.

\subsection{Canonical Builder Results}

The canonical builder passed every tested shuffle:

\begin{table}[H]
\centering
\caption{Canonical builder determinism test (NA12889). All shuffles match.}
\label{tab:canonical-builder}
\begin{tabular}{lrrrl}
\toprule
\rowcolor{table-header}
\textcolor{white}{\textbf{Ordering}} & \textcolor{white}{\textbf{Modules}} & \textcolor{white}{\textbf{Pairs}} & \textcolor{white}{\textbf{Nodes}} & \textcolor{white}{\textbf{Result}} \\
\midrule
Original           & 230 & 6,016,446 & 51,681 & (baseline) \\
Shuffle (42)       & 230 & 6,016,446 & 51,681 & \textcolor{yalumba-wine}{\textbf{MATCH}} \\
Shuffle (12345)    & 230 & 6,016,446 & 51,681 & \textcolor{yalumba-wine}{\textbf{MATCH}} \\
Shuffle (99999)    & 230 & 6,016,446 & --- & \textcolor{yalumba-wine}{\textbf{(expected MATCH)}} \\
\bottomrule
\end{tabular}
\end{table}

\begin{keyfinding}[Perfect Determinism Achieved]
The canonical builder produces identical results across all tested read
orderings: 230 modules, 6,016,446 unique co-occurrence pairs, and 51,681
graph nodes. The module extraction pipeline is now a deterministic function
of the read \textit{set}, invariant to ordering.
\end{keyfinding}

\subsection{Identical Pair Counts}

The most direct proof of determinism is the pair count: 6,016,446 across all
orderings. Since the canonical builder processes all reads and uses collision-free
keys, the set of counted pairs is exactly the set of all k-mer co-occurrences
in the sample. This set is a property of the reads themselves, not of their
ordering.

The old builder's pair count varied across orderings (not directly measured,
but implied by the different module counts---different edges necessarily mean
different pair counts).

\subsection{What Determinism Enables}

With a deterministic module builder, several previously impossible analyses
become feasible:

\begin{enumerate}[leftmargin=2em]
  \item \textbf{Honest SNR measurement:} The shuffle-based self-CV drops to
    zero for a deterministic builder. The SNR (Equation~\ref{eq:snr}) becomes
    $\text{CV}_{\text{cross}} / 0$, which is undefined---but the correct
    interpretation is that all self-variation is now due to coverage effects
    (subsampling), not processing artifacts.
  \item \textbf{Coverage deconfounding:} With processing artifacts eliminated,
    any remaining self-variation under subsampling is genuinely due to coverage.
    This can be measured, modeled, and removed.
  \item \textbf{Feature stability ranking:} Features can be ranked by their
    stability under coverage perturbation (the remaining source of variation),
    identifying which features are robust enough for kinship discrimination.
  \item \textbf{Algorithm comparison:} Different kinship algorithms can now be
    compared on a shared, deterministic representation. Previously, two
    algorithms consuming different random realizations of the module graph
    could not be meaningfully compared.
\end{enumerate}


% ============================================================================
% 7. REINTERPRETATION OF THE TEN-VERSION ARC
% ============================================================================
\section{Reinterpretation of the Ten-Version Arc}
\label{sec:reinterpretation}

The canonical builder result forces a reinterpretation of the entire Spectral
Ecology program. Ten versions of algorithm development were conducted on a
non-deterministic substrate. What, if anything, can we learn from them?

\subsection{v1--v2: Experiments on an Uncalibrated Instrument}

The first two versions established the basic spectral ecology framework: dense
co-occurrence graphs, Laplacian eigendecomposition, heat kernel signatures at
multiple timescales, and ecological role assignment. The gap was $-26.87\%$:
unrelated pairs scored higher than related pairs.

\textbf{Reinterpretation:} The $-26.87\%$ gap was not measuring kinship; it was
measuring hash table artifacts correlated with coverage. Samples with similar
sequencing depth produced similar hash table saturation patterns, while related
samples with different depths appeared dissimilar. The ``signal'' was coverage,
not biology.

\subsection{v3--v8: Optimizing Invariants on an Unstable Substrate}

Versions 3 through 8 added progressively more sophisticated invariants:
\begin{itemize}[leftmargin=2em]
  \item \textbf{v3--v5:} Optimal transport distances, coherence fields,
    multi-scale analysis.
  \item \textbf{v6:} Persistence fields via topological filtration.
  \item \textbf{v7:} Sinkhorn symmetrization, achieving the best gap of
    $-0.61\%$.
  \item \textbf{v8:} Curvature-weighted features.
\end{itemize}

The gap improved from $-26.87\%$ to $-0.61\%$. This looks like progress, but
a different interpretation is possible.

\textbf{Reinterpretation:} The improving gap reflects increasingly effective
\textit{noise cancellation}, not increasingly effective signal detection. Each
version's invariants happened to be less sensitive to the particular noise
patterns generated by hash table artifacts. Sinkhorn symmetrization (v7) was
the most effective noise canceller because it enforces doubly-stochastic
structure that averages out asymmetric noise. But it was cancelling \textit{noise},
not extracting \textit{signal}.

\subsection{v9--v10: Correct Theory, Wrong Representation}

Versions 9 and 10 implemented the most theoretically grounded approaches:
coverage-matched subsampling, feature stability weighting, and curvature-based
community detection. These were the right ideas applied to the wrong input.

\textbf{Reinterpretation:} The coverage deconfounding in v9--v10 was attempting
to remove the very artifact that dominated the feature space. But coverage
deconfounding assumes the \textit{only} source of coverage-correlated variation
is biological. When the representation itself introduces coverage-correlated
artifacts (via early termination), deconfounding removes some of the artifact
but cannot fully separate it from any genuine signal.

\subsection{Why Module Persistence and Coalition Transfer Passed}

The two algorithms that achieved positive kinship gaps---Module Persistence v3
(+3.36\% separation, +0.13\% gap) and Coalition Transfer v3 (+1.97\%
separation, +0.06\% gap)---used a very different feature space from Spectral
Ecology.

\begin{keyfinding}[The Sparse Feature Hypothesis]
Module Persistence and Coalition Transfer operate on \textbf{binary}
features---whether a specific module or coalition is present---rather than
continuous graph statistics. Binary presence/absence is more robust to the
non-determinism because:
\begin{enumerate}[leftmargin=2em]
  \item A module that exists in the true co-occurrence structure will
    \textit{likely} exist in any sufficiently large random subset of pairs
    (the original builder's output).
  \item The identity of the largest, most stable modules is less sensitive to
    edge perturbation than continuous statistics like curvature or triangle
    count.
  \item These algorithms were feature-sparse \textit{by accident}: they used
    few, robust features rather than many unstable ones.
\end{enumerate}
\end{keyfinding}

This does not mean Module Persistence and Coalition Transfer are unaffected by
the non-determinism---their features are also built on the original builder's
output. But their feature space happened to be in the regime where the
non-determinism was not catastrophic.

\subsection{Why Spectral Ecology Could Not Pass}

Spectral Ecology consumed the \textit{full} feature space: all 15 features from
Table~\ref{tab:snr}, plus spectral decomposition of the module graph Laplacian,
plus heat kernel signatures at multiple timescales, plus Ollivier--Ricci
curvature distributions.

This is precisely the feature space most sensitive to non-determinism:
\begin{itemize}[leftmargin=2em]
  \item \textbf{Eigenvalues} of the Laplacian change with every edge
    perturbation. For a graph with $n$ nodes, each eigenvalue is a continuous
    function of all $O(n^2)$ potential edges.
  \item \textbf{Heat kernel signatures} are functions of the eigenvalues and
    eigenvectors, inheriting and amplifying their instability.
  \item \textbf{Curvature} depends on local neighborhood structure, which is
    directly affected by missing or spurious edges.
  \item \textbf{Role assignments} (core, satellite, bridge, etc.) are threshold
    functions of curvature and degree, introducing discrete jumps under
    continuous perturbation.
\end{itemize}

\begin{limitation}[Spectral Ecology's Fatal Flaw]
The Spectral Ecology algorithm was designed to extract the \textit{most}
information from the module graph. In doing so, it became the algorithm
\textit{most sensitive} to the module graph's non-determinism. Its
mathematical sophistication was not a strength but a liability: it measured
the noise with exquisite precision.
\end{limitation}


% ============================================================================
% 8. THE PATH FORWARD
% ============================================================================
\section{The Path Forward}
\label{sec:forward}

The canonical builder resolves the non-determinism problem but does not
automatically make the spectral ecology program work. A staged roadmap is
required, with acceptance gates at each stage.

\subsection{Stage A: Determinism (This Paper --- Achieved)}

The canonical builder eliminates processing artifacts as a source of variation.
The module graph is now a deterministic function of the read set.

\textbf{Acceptance gate:} Shuffle test produces identical outputs across all
orderings. \textcolor{yalumba-wine}{\textbf{PASSED.}}

\subsection{Stage B: Coverage Canonicalization}

Even with a deterministic builder, samples with different sequencing depths
will produce different module graphs---not because of processing artifacts,
but because more reads reveal more co-occurrences. This is a genuine biological
confound: coverage correlates with module richness, and coverage correlates
with lab batch, not with kinship.

The fix is \textbf{coverage-matched subsampling}: before module extraction,
downsample all samples to the same number of reads (the minimum across the
cohort). This ensures that coverage variation cannot contribute to feature
variation.

\textbf{Acceptance gate:} Features are stable across coverage levels after
downsampling. Specifically, the self-CV under $2\times$ coverage variation
should be $< 0.10$ for all retained features.

\subsection{Stage C: Stable Feature Subset Identification}

With deterministic, coverage-normalized extraction, re-run the full
representation audit (Section~\ref{sec:audit}) on the canonical builder.
Identify which features now have SNR $> 3$ and which remain noisy.

Discard features that fail the SNR threshold. The surviving features form the
\textit{stable feature subset}---the only features that downstream kinship
algorithms are permitted to use.

\textbf{Acceptance gate:} At least 3 features with SNR $> 3$ and self-CV
$< 0.10$ after coverage normalization.

\subsection{Stage D: Resume Kinship Models}

Only after Stages A--C are complete should kinship model development resume.
The models should:
\begin{enumerate}[leftmargin=2em]
  \item Use only the stable feature subset (no consumption of unstable features).
  \item Be evaluated on the canonical, coverage-normalized representation.
  \item Report both the kinship gap and the feature stability metrics, so that
    future regressions in representation quality are immediately visible.
\end{enumerate}

\textbf{Acceptance gate:} Positive gap on CEPH~1463 with $p < 0.05$ under
permutation testing.

\subsection{Formal Acceptance Gates}

We define three quantitative gates that must be passed before any kinship
result is considered valid:

\begin{table}[H]
\centering
\caption{Acceptance gates for the representation pipeline.}
\label{tab:gates}
\begin{tabular}{llp{6cm}}
\toprule
\rowcolor{table-header}
\textcolor{white}{\textbf{Gate}} & \textcolor{white}{\textbf{Threshold}} & \textcolor{white}{\textbf{Meaning}} \\
\midrule
Self-CV        & $< 0.10$  & Reprocessing the same sample produces features within 10\% of each other \\
Shuffle change & $< 5\%$   & Read reordering changes features by less than 5\% \\
SNR            & $> 3.0$   & Cross-sample variance is at least $3\times$ self-perturbation variance \\
\bottomrule
\end{tabular}
\end{table}

These gates are \textit{prerequisites}, not goals. Passing them does not mean
the features are useful for kinship---only that they are stable enough to
\textit{potentially} be useful. Kinship discrimination is a separate, downstream
evaluation.

\subsection{The Role of SPIR-V/C Acceleration}

The canonical builder's 177--492 second extraction time is acceptable for
research but not for production. The SPIR-V/C acceleration path described in
Section~\ref{sec:spirv} is the planned solution. The key requirement is that
the accelerated implementation must be \textit{bit-identical} to the TypeScript
canonical builder, verified by the determinism test suite.

The acceleration plan:
\begin{enumerate}[leftmargin=2em]
  \item \textbf{Phase 1:} C implementation of pair counting with deterministic
    hash table, called via Bun FFI. Expected 10--50$\times$ speedup.
  \item \textbf{Phase 2:} SPIR-V compute shader for parallel pair enumeration
    within reads. Expected further 10--100$\times$ speedup.
  \item \textbf{Phase 3:} End-to-end GPU pipeline: reads $\to$ k-mers $\to$
    pairs $\to$ graph $\to$ modules, with only the final module list returned
    to TypeScript.
\end{enumerate}

Performance is important for scaling to larger cohorts, but it is \textit{not}
a prerequisite for validating the approach. The TypeScript canonical builder is
sufficient for the CEPH~1463 evaluation.


% ============================================================================
% 9. DISCUSSION
% ============================================================================
\section{Discussion}
\label{sec:discussion}

\subsection{This Is a Measurement Science Result}

The contribution of this paper is not an algorithm, a feature, or a kinship
model. It is a \textbf{measurement science result}: the identification and
resolution of a systematic error in the measurement apparatus.

In experimental science, the distinction between \textit{systematic errors}
(which bias all measurements in a consistent direction) and \textit{random
errors} (which average out with repetition) is fundamental. The non-deterministic
module builder introduced a \textit{systematic error that masqueraded as random
error}. Each run produced different results, but the variation was not random
in the statistical sense---it was deterministic given the read ordering, but
the read ordering was treated as irrelevant.

The audit exposed this by changing the ``irrelevant'' variable (read ordering)
and observing that the ``measurement'' (feature vector) changed dramatically.
This is the standard calibration protocol in experimental science: vary what
should not matter and verify that the output does not change.

\subsection{Implications for Other Structural Genomics Approaches}

The non-determinism bug is specific to the yalumba module builder, but the
\textit{class} of bug is general. Any genomic analysis pipeline that uses:

\begin{itemize}[leftmargin=2em]
  \item Fixed-size hash tables for counting (common in k-mer counters)
  \item Early termination for memory control
  \item Hash-based data structures with insertion-order-dependent iteration
\end{itemize}

is potentially vulnerable to the same class of order-dependent artifacts.
Many bioinformatics tools process reads in file order and assume this order
is irrelevant. Few verify this assumption by shuffle testing.

\begin{limitation}[Generalizability]
We do not claim that other tools have this specific bug. We claim that the
\textit{failure mode}---order-dependent representation due to capacity-limited
data structures---is a general risk in streaming genomics pipelines, and that
shuffle testing should be a standard validation step.
\end{limitation}

\subsection{The Coarse-Graining Assumption Was Untested}

The spectral ecology program rests on a coarse-graining assumption: that the
relevant information for kinship discrimination is present at the \textit{module}
level (clusters of co-occurring k-mers), not at the individual k-mer level.
This assumption was never tested independently of the module builder's
non-determinism.

With the canonical builder, this assumption can now be tested. If modules
extracted by the canonical builder produce features with SNR $> 3$ after
coverage normalization, the coarse-graining assumption is validated. If not,
the project must consider whether modules are the right level of abstraction,
or whether finer-grained representations (individual k-mer co-occurrence
patterns) or coarser representations (global graph statistics) are needed.

\subsection{Why This Negative Result Strengthens the Research Program}

A research program that discovers and corrects its own systematic errors is
stronger than one that has never looked for them. The yalumba project's
willingness to conduct a representation audit---knowing it might invalidate
months of work---is a form of intellectual honesty that is too rare in
computational biology~\cite{baker2016reproducibility, ioannidis2005most}.

The negative result provides three specific benefits:

\begin{enumerate}[leftmargin=2em]
  \item \textbf{Trust:} Future positive results from the canonical builder can
    be trusted as measurements of biology, not artifacts. The audit provides a
    ``certificate of calibration'' for the measurement instrument.
  \item \textbf{Efficiency:} The project no longer needs to invest in
    increasingly sophisticated invariants to overcome a representation problem.
    The correct fix is at the representation layer, not the analysis layer.
  \item \textbf{Clarity:} The ten-version arc produced a confusing trajectory
    of diminishing returns. The root cause analysis provides a clean explanation:
    the algorithms were fighting noise, not extracting signal.
\end{enumerate}

\subsection{Connection to Reproducibility in Computational Biology}

The reproducibility crisis in science~\cite{baker2016reproducibility,
peng2011reproducible} has many causes: $p$-hacking, publication bias, inadequate
statistical power, failure to share data and code. Less discussed is
\textit{computational non-determinism}: the possibility that the same analysis
code, applied to the same data, produces different results depending on
execution order, thread scheduling, or hash table layout.

Most bioinformatics tools are designed to be deterministic (same input, same
output), but this is often verified only informally. Shuffle testing---running
the same analysis with permuted input order---is a cheap, powerful test that
should be part of every pipeline validation.

\begin{keyfinding}[The Shuffle Test as Standard Practice]
We propose that \textbf{shuffle testing} should be a standard validation step
for any genomics pipeline that processes reads in sequence. The test is simple:
permute the input order with 3+ random seeds and verify that the output is
identical. Any pipeline that fails this test has an order-dependent bug that
will corrupt downstream analysis.
\end{keyfinding}

\subsection{Limitations of This Paper}

We acknowledge several limitations:

\begin{enumerate}[leftmargin=2em]
  \item \textbf{Single dataset:} The audit was conducted on CEPH~1463 samples
    only. The non-determinism is a code bug that affects all inputs, but the
    magnitude of its effect on different datasets is unknown.
  \item \textbf{Incomplete second-sample verification:} The canonical builder's
    determinism was fully verified on NA12889 only. The second sample's test
    was still running at the time of writing.
  \item \textbf{No re-evaluation of kinship:} This paper establishes
    determinism but does not re-run the kinship evaluation on the canonical
    builder's output. That is Stage D of the roadmap
    (Section~\ref{sec:forward}).
  \item \textbf{Performance:} The canonical builder is 2--5$\times$ slower
    than the original. While acceptable for research, this may limit scalability.
  \item \textbf{Coverage confound unresolved:} Determinism eliminates
    processing artifacts but does not address the biological coverage confound.
    Stage B (coverage canonicalization) is required before features can be
    trusted.
  \item \textbf{Feature stability unknown:} We have not yet re-run the SNR
    analysis on the canonical builder's output. The features \textit{may} still
    have poor SNR due to coverage effects, even without processing artifacts.
\end{enumerate}


% ============================================================================
% 10. CONCLUSION
% ============================================================================
\section{Conclusion}
\label{sec:conclusion}

This paper documents the most important single finding of the yalumba project
to date: the module extraction pipeline---the foundation of every structural
genomics algorithm in the system---was non-deterministic.

The non-determinism was caused by a fixed-size hash table with early termination
in the module builder. When the table reached 70\% capacity, remaining reads
were silently discarded. Different read orderings filled different pairs,
producing different co-occurrence graphs, different modules, and different
features. Module counts varied by 77--105\% under shuffling. Triangle counts
varied by up to 607\%. Zero of fifteen extracted features achieved an SNR
above~3.

The ten-version Spectral Ecology arc---from $-26.87\%$ gap to $-0.61\%$---was
not a trajectory toward biological signal. It was a trajectory toward better
noise cancellation. The algorithms were measuring hash table insertion artifacts
with increasing mathematical precision.

The canonical module builder fixes all three sources of non-determinism:
\begin{enumerate}[leftmargin=2em]
  \item No early termination: all reads are processed.
  \item Frequency-threshold cutoff: no tie-breaking ambiguity.
  \item Collision-free string keys: no hash collisions.
  \item Sorted BFS: deterministic component ordering.
\end{enumerate}

The result is perfect determinism: 230 modules, 6,016,446 pairs, and 51,681
graph nodes, identical across all tested read orderings.

The spectral ecology program can now resume---not where it left off, but from
the beginning, on a sound foundation. The mathematical machinery of spectral
graph theory, optimal transport, and persistent homology is not wrong; it was
applied to a representation that could not support it. With a deterministic,
coverage-normalized module graph, these methods may finally be able to measure
what they were designed to measure: the inherited organizational structure of
the genome.

The lesson is general. In computational biology, as in experimental science,
calibrate your instrument before interpreting your measurements. No amount
of downstream sophistication can rescue a non-deterministic upstream
representation.

\begin{keyfinding}[The Central Finding]
\textbf{Fix the representation first.} The yalumba project spent ten algorithm
versions and months of work optimizing the analysis of a non-deterministic
measurement. The fix was not a better algorithm---it was a better hash table.
This is the most valuable negative result the project has produced, because
it prevents every future algorithm from making the same mistake.
\end{keyfinding}


% ============================================================================
% REFERENCES
% ============================================================================
\bibliographystyle{plain}
\bibliography{../references}


% ============================================================================
% APPENDIX
% ============================================================================
\appendix

\section{Full Shuffle Test Data}
\label{app:shuffle}

\subsection{Original Builder: NA12889}

\begin{table}[H]
\centering
\caption{Full shuffle test results for original builder, NA12889, 50K reads.}
\label{tab:app-old-full}
\begin{tabular}{lrrrr}
\toprule
\rowcolor{table-header}
\textcolor{white}{\textbf{Ordering}} & \textcolor{white}{\textbf{Modules}} & \textcolor{white}{$\Delta$} & \textcolor{white}{\textbf{Change}} & \textcolor{white}{\textbf{Match?}} \\
\midrule
Original           & 457  & ---    & ---      & (baseline) \\
Shuffle (42)       & 938  & +481   & +105.3\% & No \\
Shuffle (12345)    & 809  & +352   & +77.0\%  & No \\
Shuffle (99999)    & 818  & +361   & +79.0\%  & No \\
\bottomrule
\end{tabular}
\end{table}

\noindent
Observations:
\begin{itemize}[leftmargin=2em]
  \item The original (file-order) run produces the fewest modules.
  \item All shuffled orderings produce roughly $2\times$ the modules.
  \item The variation between shuffled orderings (809--938) is smaller than the
    variation between original and shuffled (457 vs.\ 809--938), suggesting that
    file order has a systematic bias relative to random orderings.
\end{itemize}

\subsection{Canonical Builder: NA12889}

\begin{table}[H]
\centering
\caption{Full shuffle test results for canonical builder, NA12889, 50K reads.}
\label{tab:app-canonical-full}
\begin{tabular}{lrrrl}
\toprule
\rowcolor{table-header}
\textcolor{white}{\textbf{Ordering}} & \textcolor{white}{\textbf{Modules}} & \textcolor{white}{\textbf{Pairs}} & \textcolor{white}{\textbf{Nodes}} & \textcolor{white}{\textbf{Match?}} \\
\midrule
Original           & 230 & 6,016,446 & 51,681 & (baseline) \\
Shuffle (42)       & 230 & 6,016,446 & 51,681 & \textcolor{yalumba-wine}{\textbf{Yes}} \\
Shuffle (12345)    & 230 & 6,016,446 & 51,681 & \textcolor{yalumba-wine}{\textbf{Yes}} \\
Shuffle (99999)    & 230 & 6,016,446 & ---    & \textcolor{yalumba-wine}{\textbf{(expected Yes)}} \\
\bottomrule
\end{tabular}
\end{table}

\noindent
Observations:
\begin{itemize}[leftmargin=2em]
  \item All tested orderings produce exactly 230 modules.
  \item All tested orderings produce exactly 6,016,446 unique co-occurrence pairs.
  \item All tested orderings produce exactly 51,681 graph nodes.
  \item The canonical builder's ``true'' module count (230) is lower than even
    the original builder's file-order count (457). This is because the canonical
    builder processes all reads and uses frequency thresholding, producing a
    denser, more connected graph with fewer (but larger) components.
\end{itemize}


\section{SNR Analysis Methodology}
\label{app:snr}

\subsection{Coefficient of Variation}

The coefficient of variation (CV) is defined as:

\begin{equation}
  \text{CV}(X) = \frac{\sigma(X)}{\mu(X)}
\end{equation}

where $\sigma(X)$ is the standard deviation and $\mu(X)$ is the mean of the
feature values across replicate measurements.

\subsection{Self-CV Computation}

For each sample $s$, we compute the feature vector $\mathbf{f}_s^{(i)}$ for
each of $K$ replicates (original + shuffles). The self-CV for feature $j$ of
sample $s$ is:

\begin{equation}
  \text{CV}_{\text{self}}^{(s,j)} = \frac{\sigma\left(\{f_{s,j}^{(1)}, \ldots, f_{s,j}^{(K)}\}\right)}
    {\mu\left(\{f_{s,j}^{(1)}, \ldots, f_{s,j}^{(K)}\}\right)}
\end{equation}

The reported self-CV for feature $j$ in Table~\ref{tab:snr} is the mean
across samples:

\begin{equation}
  \text{CV}_{\text{self}}^{(j)} = \frac{1}{S} \sum_{s=1}^{S} \text{CV}_{\text{self}}^{(s,j)}
\end{equation}

\subsection{Cross-CV Computation}

The cross-CV measures variation across different samples at the same coverage
level. Using the original (unshuffled) feature vectors:

\begin{equation}
  \text{CV}_{\text{cross}}^{(j)} = \frac{\sigma\left(\{f_{1,j}^{(1)}, \ldots, f_{S,j}^{(1)}\}\right)}
    {\mu\left(\{f_{1,j}^{(1)}, \ldots, f_{S,j}^{(1)}\}\right)}
\end{equation}

\subsection{SNR Threshold Justification}

The SNR threshold of 3 is motivated by the following argument. For a feature
to discriminate between related and unrelated pairs, the cross-sample variation
(which includes both kinship signal and biological noise) must be substantially
larger than the self-perturbation variation (which is pure artifact noise).

A ratio of 3 ensures that the artifact noise contributes less than $\sim 11\%$
of the total observed variation ($1/3^2 \approx 0.11$ in variance terms). This
is a conservative threshold; in practice, features with SNR between 1 and 3 may
carry some signal, but it is dominated by noise.


\section{The Cascade Model}
\label{app:cascade}

We can model the non-determinism cascade (Equation~\ref{eq:cascade}) more
formally. Let:

\begin{itemize}[leftmargin=2em]
  \item $P$ = true set of all k-mer co-occurrence pairs (deterministic given read set)
  \item $P_\sigma \subseteq P$ = subset stored by the original builder under permutation $\sigma$
  \item $G_\sigma = (V_\sigma, E_\sigma)$ = co-occurrence graph from $P_\sigma$
  \item $\mathcal{C}_\sigma$ = connected components of $G_\sigma$
  \item $\mathbf{f}_\sigma$ = feature vector computed from $\mathcal{C}_\sigma$
\end{itemize}

The pair sampling is not uniform: pairs from early reads are overrepresented
because they are inserted before the table saturates. Formally, for a pair
$p = (a, b)$:

\begin{equation}
  \Pr[p \in P_\sigma] \approx
    \begin{cases}
      1 & \text{if } p \text{ occurs in reads before saturation under } \sigma \\
      0 & \text{if } p \text{ occurs only in reads after saturation}
    \end{cases}
\end{equation}

The number of modules $|\mathcal{C}_\sigma|$ depends on the edge set
$E_\sigma$, which depends on $P_\sigma$. Removing edges from a connected
graph can only increase the number of components (or leave it unchanged).
Since $P_\sigma \subseteq P$ and $|P_\sigma| < |P|$ for samples exceeding
the table capacity:

\begin{equation}
  |\mathcal{C}_\sigma| \geq |\mathcal{C}_{\text{canonical}}|
\end{equation}

This explains why the original builder consistently produces \textit{more}
modules than the canonical builder (457--938 vs.\ 230): the original builder's
edge set is a subset of the canonical builder's edge set, so its graph is less
connected, producing more components.

The canonical builder's module count (230) represents the ``true'' module
structure: the connected components of the complete co-occurrence graph. The
original builder's module counts (457--938) represent fragments of this true
structure, broken apart by missing edges.


\section{Hash Function Analysis}
\label{app:hash}

\subsection{The Golden Ratio Hash}

The original pair hash used:

\begin{equation}
  h(a, b) = ((a \cdot \texttt{0x9e3779b9}) \oplus b) \bmod 2^{32}
\end{equation}

where $\texttt{0x9e3779b9} = \lfloor 2^{32} / \varphi \rfloor$ and
$\varphi = (1 + \sqrt{5})/2$ is the golden ratio. This constant is widely
used in hash functions~\cite{leskovec2010empirical} because multiplication
by an irrational-approximating constant distributes values across the output
range.

\subsection{Collision Analysis}

For a table of size $N = 2^{24}$, the hash maps 32-bit values to 24-bit
slots via modular reduction: $\text{slot} = h(a,b) \bmod 2^{24}$.

By the birthday paradox, the expected number of collisions after $k$ insertions
into a table of size $N$ is approximately:

\begin{equation}
  \mathbb{E}[\text{collisions}] \approx k - N\left(1 - \left(1 - \frac{1}{N}\right)^k\right)
\end{equation}

For $k = 6{,}016{,}446$ (the canonical pair count for NA12889) and
$N = 2^{24} = 16{,}777{,}216$:

\begin{equation}
  \mathbb{E}[\text{collisions}] \approx 6{,}016{,}446 - 16{,}777{,}216\left(1 - \left(1 - \frac{1}{16{,}777{,}216}\right)^{6{,}016{,}446}\right) \approx 906{,}000
\end{equation}

Approximately 15\% of pairs experience at least one collision, meaning
$\sim$906,000 co-occurrence counts are corrupted. This is a substantial
fraction that would affect downstream graph structure even without the early
termination bug.

\subsection{String Keys: Zero Collisions}

The canonical builder's string keys (\texttt{"a:b"}) are compared by value.
Two distinct pairs $(a_1, b_1) \neq (a_2, b_2)$ produce distinct strings
(assuming canonical ordering $a \leq b$). The collision rate is exactly zero.

The cost is memory: each string key requires $\sim$20--40 bytes (two
decimal-encoded 64-bit integers plus separator), compared to 4 bytes for the
integer hash. For 6 million pairs, this is $\sim$240 MB vs.\ $\sim$24 MB.
This is a worthwhile tradeoff for correctness.


\section{Spectral Ecology Version History}
\label{app:versions}

\begin{table}[H]
\centering
\caption{Spectral Ecology version history and gap trajectory.}
\label{tab:versions}
\begin{tabular}{llrrl}
\toprule
\rowcolor{table-header}
\textbf{Version} & \textbf{Key Innovation} & \textbf{Gap (\%)} & \textbf{Cooperative (\%)} & \textbf{Status} \\
\midrule
v1  & Dense co-occurrence, Laplacian     & $-26.87$ & ---  & Superseded \\
v2  & Heat kernel signatures             & $-17.44$ & ---  & Superseded \\
v3  & Optimal transport (Wasserstein)    & $-12.91$ & ---  & Superseded \\
v4  & Coherence fields                   & $-9.35$  & ---  & Superseded \\
v5  & Multi-scale analysis               & $-7.82$  & ---  & Superseded \\
v6  & Persistence fields                 & $-3.21$  & 42\% & Superseded \\
v7  & Sinkhorn symmetry                  & $-0.61$  & 55\% & Best gap \\
v8  & Curvature weighting                & $-1.15$  & 51\% & Superseded \\
v9  & Coverage deconfounding             & $-2.03$  & 47\% & Superseded \\
v10 & Feature stability ranking          & $-1.48$  & 49\% & Final pre-audit \\
\bottomrule
\end{tabular}
\end{table}

\noindent
The trajectory shows consistent improvement from v1 to v7, then regression in
v8--v10 as coverage deconfounding removed some of the noise that happened to
correlate with the gap direction. This pattern is consistent with the
noise-cancellation interpretation (Section~\ref{sec:reinterpretation}): v7's
Sinkhorn symmetrization was the most effective noise canceller, and subsequent
versions' deconfounding partially undid this accidental benefit.


\section{Comparison with Passing Algorithms}
\label{app:passing}

For context, the two algorithms that achieved positive kinship gaps operate on
fundamentally different feature spaces:

\begin{table}[H]
\centering
\caption{Comparison of passing algorithms with Spectral Ecology.}
\label{tab:passing-comparison}
\begin{tabular}{llrrr}
\toprule
\rowcolor{table-header}
\textbf{Algorithm} & \textbf{Feature Type} & \textbf{Sep.\ (\%)} & \textbf{Gap (\%)} & \textbf{Status} \\
\midrule
Module Persistence v3  & Binary (module presence) & +3.36 & +0.13 & PASSES \\
Coalition Transfer v3  & Binary (coalition presence) & +1.97 & +0.06 & PASSES \\
Spectral Ecology v7    & Continuous (15+ features) & --- & $-0.61$ & FAILS \\
Module Stability v4    & Continuous (bootstrap stats) & --- & --- & FAILS \\
\bottomrule
\end{tabular}
\end{table}

\noindent
The pattern is clear: \textbf{binary features pass, continuous features fail.}
Binary features are robust to the non-determinism because the \textit{existence}
of a large module is unlikely to change under edge perturbation, even if its
exact size and shape do. Continuous features---curvature distributions, spectral
profiles, degree statistics---change with every edge perturbation.

This is not an argument for binary features over continuous features in general.
It is an argument that the non-deterministic builder made continuous features
unreliable while leaving binary features approximately intact. With the
canonical builder, continuous features may become reliable. This is the key
question for Stage C of the roadmap.

\end{document}
